\section{Conclusion}
\label{sec:conclusion}

We introduced Adaptive Dual-Branch Decoding, a training-free method for reducing hallucinations in LVLMs. ADBD replaces ONLY's frozen one-layer binary decision with a continuous TVER state refined across selected decoder layers by head-specific adaptive gains. The state produces complementary high- and low-TVER auxiliary branches, which are independently reinforced or penalized according to their disagreement with the normal distribution. An anchored three-branch rule combines these signals without requiring a separate full LVLM invocation.

On the full POPE benchmark with LLaVA-1.5-7B, ADBD achieves 82.10\%, 86.37\%, and 89.13\% accuracy on Random, Popular, and Adversarial, respectively, yielding 85.87\% macro accuracy and 85.48\% macro F1. Precision increases across the three splits while recall remains stable, producing the strongest result on Adversarial negatives. On MME Hallucination, the method improves Position and Color but decreases Existence and Count, yielding a reported aggregate score of 620.00 compared with 635.55 for ONLY; the aggregate must be verified because the supplied ADBD category scores sum to 630.00. Overall, the results suggest that complementary and adaptively accumulated attention evidence can provide a useful contrastive signal, but its gains are category-dependent rather than universal.
